% PAGES: 141-142
% Continues sec04_c2_3.tex (pages 134-140 / folios 118-124), whose last
% line was the completed statement of Theorem 4.6 (Gershgorin's Theorem)
% followed by a grammatically complete explanatory remark -- no environment
% was left open across that boundary. Per c2's own note, the "Proof." of
% Theorem 4.6 begins here, on page 141 (folio 125), as a fresh
% \begin{proof}...\end{proof} -- it is NOT a continuation of anything
% opened in sec04_c2_3.tex.
%
% Section 4.3 "Diagonally dominant matrices" continues: proof of Theorem
% 4.6, then Theorem 4.7 (with proof), Definition 4.9, and Lemma 4.9 (with
% proof). Both pages close cleanly -- Lemma 4.9's proof ends with the
% printed square at the bottom of page 142 (folio 126). No environment is
% left open at the end of this file.

\begin{proof}
Let $\lambda$ be an eigenvalue of $A$, and let $v$ be a corresponding eigenvector. Let
$j$ be the index of the largest entry of $v$ in absolute value, i.e., such that
\begin{equation}
|v_j| \;=\; \max_{k=1:n} |v_k|.
\end{equation}
Then,
\begin{equation}
|v_k| \;\leq\; |v_j|, \quad \forall\, k = 1:n.
\end{equation}
Since $Av = \lambda v$, it follows that
\begin{equation}
\sum_{k=1}^n A(i,k) v_k \;=\; \lambda v_i, \quad \forall\, i = 1:n.
\end{equation}
By letting $i = j$ in (4.34), where $j$ is the index of the largest entry of $v$ in
absolute value, see (4.32), we find that
\[
\lambda v_j \;=\; \sum_{k=1}^n A(j,k)v_k \;=\; A(j,j)v_j \;+\; \sum_{k=1:n,\, k \neq j}
A(j,k)v_k,
\]
and therefore
\begin{equation}
(\lambda - A(j,j))v_j \;=\; \sum_{k=1:n,\, k \neq j} A(j,k)v_k.
\end{equation}
By taking absolute values in (4.35) and dividing by $|v_j|$, we obtain, using the
notation from (4.28), that
\begin{align*}
|\lambda - A(j,j)| \;=\; \frac{\left|\sum_{k=1:n,\, k \neq j} A(j,k)v_k\right|}{|v_j|}
&\;\leq\; \frac{\sum_{k=1:n,\, k \neq j} |A(j,k)||v_k|}{|v_j|}\\
&\;=\; \sum_{k=1:n,\, k \neq j} |A(j,k)| \frac{|v_k|}{|v_j|}\\
&\;\leq\; \sum_{k=1:n,\, k \neq j} |A(j,k)|\\
&\;=\; R_j,
\end{align*}
since $\frac{|v_k|}{|v_j|} \leq 1$; cf. (4.33).
\end{proof}

Gershgorin's Theorem can be used to derive useful properties of diagonally dominant
matrices, see Theorem 4.7 and Theorem 5.7.

\begin{theorem}
Any strictly diagonally dominant matrix is nonsingular.
\end{theorem}

\begin{proof}
Let $\lambda$ be an eigenvalue of $A$. From Gershgorin's Theorem, it follows that there
exists $j$, with $1 \leq j \leq n$, such that
\begin{equation}
|A(j,j) - \lambda| \;\leq\; R_j;
\end{equation}
cf. (4.31), and using the fact that $|\lambda - A(j,j)| = |A(j,j) - \lambda|$. Note that
$|a| - |b| \leq |a-b|$, for any $a, b \in \C$. Then,
\begin{equation}
|A(j,j)| - |\lambda| \;\leq\; |A(j,j) - \lambda|.
\end{equation}
From (4.36) and (4.37), it follows that
\[
|A(j,j)| - |\lambda| \;\leq\; R_j,
\]
and therefore
\[
|\lambda| \;\geq\; |A(j,j)| - R_j \;>\; 0,
\]
since $A$ is a strictly diagonally dominant matrix; cf. (4.30).
Thus, $|\lambda| > 0$, and therefore $\lambda \neq 0$. We conclude that any eigenvalue
of the matrix $A$ is nonzero, and therefore that the matrix $A$ is nonsingular; cf.
Theorem 4.3.
\end{proof}

Note that, although less frequently used in practice, column--based versions of
Gershgorin's Theorem and of Theorem 4.7 also hold, and are based on the fact that a
square matrix and its transpose have the same eigenvalues; see Lemma 4.5. We only
provide a version of Theorem 4.7 here; see Theorem 4.8.

\begin{definition}
The $n \times n$ matrix $A$ is strictly column diagonally dominant if and only if
\begin{equation}
|A(k,k)| \;>\; \sum_{j=1:n,\, j \neq k} |A(j,k)|, \quad \forall\, k = 1:n,
\end{equation}
i.e., if and only if, for every column, the absolute value of the main diagonal entry
from the column is strictly greater than the sum of the absolute values of all the
other entries in that column.
\end{definition}

Note that, unless otherwise specified, diagonal dominance will refer to the ``row''
diagonal dominance introduced in Definitions 4.7 and 4.8.

\begin{lemma}
If $A$ is a strictly column diagonally dominant matrix, then its transpose $A^t$ is a
strictly diagonally dominant matrix.
\end{lemma}

\begin{proof}
Let $A$ be a strictly column diagonally dominant matrix. From (4.38), it follows that
\begin{equation}
|A(k,k)| \;>\; \sum_{j=1:n,\, j \neq k} |A(j,k)|,
\end{equation}
for all $k = 1:n$. By switching the indices $j$ and $k$ in the inequality (4.38), we
find that
\begin{equation}
|A(j,j)| \;>\; \sum_{k=1:n,\, k \neq j} |A(k,j)|,
\end{equation}
for all $j = 1:n$. Note that
\begin{equation}
\sum_{k=1:n,\, k \neq j} |A(k,j)| \;=\; \sum_{k=1:n,\, k \neq j} |A^t(j,k)|,
\end{equation}
since, by definition, $A^t(j,k) = A(k,j)$ for all $1 \leq j,k \leq n$. Since
$A(j,j) = A^t(j,j)$, it follows from (4.40) and (4.41) that
\[
|A^t(j,j)| \;>\; \sum_{k=1:n,\, k \neq j} |A^t(j,k)|.
\]
Then, from Definition 4.8, we find that the matrix $A^t$ is strictly diagonally
dominant.
\end{proof}
